% GENERATED FILE. Edit canonical lessons, then run npm run generate:books.
% canonical-source-hash: fnv1a64:325837d0a3e40cb2
% canonical-lessons: ML-C393-anusmaranam, ML-C393-smarakam, ML-C393-pratima, ML-C393-silalikhitam, ML-C393-anusmarikkuka

\chapter{Monuments and Memory}
\label{ch:ml-monuments-and-memory}
\hlchaptermodality{393}

\begin{chapteropening}
\textbf{By the end of this chapter:} \emph{I can say a commemoration, a memorial, a statue, a stone inscription and to commemorate.}

Complete the last lesson of chapter 393: I can say a commemoration, a memorial, a statue, a stone inscription and to commemorate.
\end{chapteropening}

\section[\texorpdfstring{\ml{അനുസ്മരണം} (anusmaraṇaṁ)}{അനുസ്മരണം (anusmaraṇaṁ)}]{\ml{അനുസ്മരണം} (anusmaraṇaṁ) — a commemoration}

\label{lesson:ML-C393-anusmaranam}



\begin{quote}
Before the new one: say the Malayalam for a cemetery, then the Malayalam for to offer condolences.
\end{quote}

\subsection*{You'll want to know: \ml{അനുസ്മരണം}}
\textbf{\ml{അനുസ്മരണം}} — \emph{anusmaraṇaṁ} — \textquotedblleft{}a commemoration\textquotedblright{}.

\begin{quote}
\ml{കവിയുടെ} \ml{അനുസ്മരണം} \ml{ഗ്രന്ഥശാലയിൽ} \ml{നടന്നു}.
\end{quote}

\begin{quote}
\emph{kaviyuṭe anusmaraṇaṁ granthaśālayil naṭannu.}
\end{quote}

\begin{quote}
The poet's commemoration took place in the library.
\end{quote}

A gathering held to remember someone who has died.

\begin{cousinweb}
\textbf{\ml{അനുസ്മരണം}} (\emph{anusmaraṇaṁ}) is borrowed from Sanskrit \emph{anusmaraṇa}, \textquotedblleft{}remembering, repeated recollection\textquotedblright{} (Monier-Williams).
\end{cousinweb}

\begin{grammarlens}[title={What you've built}]
One more word for telling what happened.
\end{grammarlens}

\subsection*{Your turn}
\begin{itemize}
\raggedright
  \item \emph{Say it:} \emph{anusmaraṇaṁ}
  \item \emph{Say it:} \emph{anusmaraṇaṁ}, once more
  \item \emph{Say it:} \emph{kaviyuṭe anusmaraṇaṁ granthaśālayil naṭannu.} — The poet's commemoration took place in the library.
  \item \emph{Recall:} say the Malayalam for a cemetery, then the Malayalam for to offer condolences, then say \emph{anusmaraṇaṁ} again
\end{itemize}

\begin{tcolorbox}[breakable,colback=teal!4,colframe=teal!35!black,title={Before you move on}]
What does \ml{അനുസ്മരണം} mean? (\textquotedblleft{}A commemoration\textquotedblright{}.) Say it once more.
\end{tcolorbox}
\section[\texorpdfstring{\ml{സ്മാരകം} (smārakaṁ)}{സ്മാരകം (smārakaṁ)}]{\ml{സ്മാരകം} (smārakaṁ) — a memorial; a monument}

\label{lesson:ML-C393-smarakam}



\begin{quote}
Before the new one: say the Malayalam for to offer condolences, then the Malayalam for a commemoration.
\end{quote}

\subsection*{You'll want to know: \ml{സ്മാരകം}}
\textbf{\ml{സ്മാരകം}} — \emph{smārakaṁ} — \textquotedblleft{}a memorial; a monument\textquotedblright{}.

\begin{quote}
\ml{യുദ്ധത്തിൽ} \ml{മരിച്ചവരുടെ} \ml{സ്മാരകം} \ml{നഗരത്തിന്റെ} \ml{നടുവിൽ} \ml{ആണ്}.
\end{quote}

\begin{quote}
\emph{yuddhattil mariccavaruṭe smārakaṁ nagarattinṟe naṭuvil āṇŭ.}
\end{quote}

\begin{quote}
The memorial to those who died in the war is in the middle of the city.
\end{quote}

Anything built to keep a memory: a statue, a hall, a pillar.

\begin{cousinweb}
\textbf{\ml{സ്മാരകം}} (\emph{smārakaṁ}) is borrowed from Sanskrit \emph{smāraka}, \textquotedblleft{}recalling, reminding of\textquotedblright{} (Monier-Williams).
\end{cousinweb}

\begin{grammarlens}[title={What you've built}]
One more word for telling what happened.
\end{grammarlens}

\subsection*{Your turn}
\begin{itemize}
\raggedright
  \item \emph{Say it:} \emph{smārakaṁ}
  \item \emph{Say it:} \emph{smārakaṁ}, once more
  \item \emph{Say it:} \emph{yuddhattil mariccavaruṭe smārakaṁ nagarattinṟe naṭuvil āṇŭ.} — The memorial to those who died in the war is in the middle of the city.
  \item \emph{Recall:} say the Malayalam for to offer condolences, then the Malayalam for a commemoration, then say \emph{smārakaṁ} again
\end{itemize}

\begin{tcolorbox}[breakable,colback=teal!4,colframe=teal!35!black,title={Before you move on}]
What does \ml{സ്മാരകം} mean? (\textquotedblleft{}A memorial; a monument\textquotedblright{}.) Say it once more.
\end{tcolorbox}
\section[\texorpdfstring{\ml{പ്രതിമ} (pratima)}{പ്രതിമ (pratima)}]{\ml{പ്രതിമ} (pratima) — a statue}

\label{lesson:ML-C393-pratima}



\begin{quote}
Before the new one: say the Malayalam for a commemoration, then the Malayalam for a memorial.
\end{quote}

\subsection*{You'll want to know: \ml{പ്രതിമ}}
\textbf{\ml{പ്രതിമ}} — \emph{pratima} — \textquotedblleft{}a statue\textquotedblright{}.

\begin{quote}
\ml{കവലയിൽ} \ml{ഒരു} \ml{കവിയുടെ} \ml{വലിയ} \ml{പ്രതിമ} \ml{ഉണ്ട്}.
\end{quote}

\begin{quote}
\emph{kavalayil oru kaviyuṭe valiya pratima uṇṭŭ.}
\end{quote}

\begin{quote}
There is a big statue of a poet at the junction.
\end{quote}

Plural \ml{പ്രതിമകൾ}.

\begin{cousinweb}
\textbf{\ml{പ്രതിമ}} (\emph{pratima}) is borrowed from Sanskrit \emph{pratimā}, \textquotedblleft{}an image, a likeness\textquotedblright{} (CDIAL 8585).
\end{cousinweb}

\begin{grammarlens}[title={What you've built}]
One more word for telling what happened.
\end{grammarlens}

\subsection*{Your turn}
\begin{itemize}
\raggedright
  \item \emph{Say it:} \emph{pratima}
  \item \emph{Say it:} \emph{pratima}, once more
  \item \emph{Say it:} \emph{kavalayil oru kaviyuṭe valiya pratima uṇṭŭ.} — There is a big statue of a poet at the junction.
  \item \emph{Recall:} say the Malayalam for a commemoration, then the Malayalam for a memorial, then say \emph{pratima} again
\end{itemize}

\begin{tcolorbox}[breakable,colback=teal!4,colframe=teal!35!black,title={Before you move on}]
What does \ml{പ്രതിമ} mean? (\textquotedblleft{}A statue\textquotedblright{}.) Say it once more.
\end{tcolorbox}
\section[\texorpdfstring{\ml{ശിലാലിഖിതം} (śilālikhitaṁ)}{ശിലാലിഖിതം (śilālikhitaṁ)}]{\ml{ശിലാലിഖിതം} (śilālikhitaṁ) — a stone inscription}

\label{lesson:ML-C393-silalikhitam}



\begin{quote}
Before the new one: say the Malayalam for a memorial, then the Malayalam for a statue.
\end{quote}

\subsection*{You'll want to know: \ml{ശിലാലിഖിതം}}
\textbf{\ml{ശിലാലിഖിതം}} — \emph{śilālikhitaṁ} — \textquotedblleft{}a stone inscription\textquotedblright{}.

\begin{quote}
\ml{ക്ഷേത്രത്തിന്റെ} \ml{ഭിത്തിയിൽ} \ml{ഒരു} \ml{പഴയ} \ml{ശിലാലിഖിതം} \ml{ഉണ്ട്}.
\end{quote}

\begin{quote}
\emph{kṣētrattinṟe bhittiyil oru paḻaya śilālikhitaṁ uṇṭŭ.}
\end{quote}

\begin{quote}
There is an old stone inscription on the temple's wall.
\end{quote}

\ml{ശില}, stone, + \ml{ലിഖിതം}, something written.

\begin{cousinweb}
\textbf{\ml{ശിലാലിഖിതം}} (\emph{śilālikhitaṁ}) is built from two Sanskrit words, \emph{śilā}, \textquotedblleft{}a rock, a crag\textquotedblright{} (CDIAL 12459), and \emph{likhita}, \textquotedblleft{}scraped; written\textquotedblright{} (CDIAL 11050): writing cut into stone.
\end{cousinweb}

\begin{grammarlens}[title={What you've built}]
One more word for telling what happened.
\end{grammarlens}

\subsection*{Your turn}
\begin{itemize}
\raggedright
  \item \emph{Say it:} \emph{śilālikhitaṁ}
  \item \emph{Say it:} \emph{śilālikhitaṁ}, once more
  \item \emph{Say it:} \emph{kṣētrattinṟe bhittiyil oru paḻaya śilālikhitaṁ uṇṭŭ.} — There is an old stone inscription on the temple's wall.
  \item \emph{Recall:} say the Malayalam for a memorial, then the Malayalam for a statue, then say \emph{śilālikhitaṁ} again
\end{itemize}

\begin{tcolorbox}[breakable,colback=teal!4,colframe=teal!35!black,title={Before you move on}]
What does \ml{ശിലാലിഖിതം} mean? (\textquotedblleft{}A stone inscription\textquotedblright{}.) Say it once more.
\end{tcolorbox}
\section[\texorpdfstring{\ml{അനുസ്മരിക്കുക} (anusmarikkuka)}{അനുസ്മരിക്കുക (anusmarikkuka)}]{\ml{അനുസ്മരിക്കുക} (anusmarikkuka) — to commemorate}

\label{lesson:ML-C393-anusmarikkuka}



\begin{quote}
Before the new one: say the Malayalam for a statue, then the Malayalam for a stone inscription.
\end{quote}

\subsection*{You'll want to know: \ml{അനുസ്മരിക്കുക}}
\textbf{\ml{അനുസ്മരിക്കുക}} — \emph{anusmarikkuka} — \textquotedblleft{}to commemorate\textquotedblright{}.

\begin{quote}
\ml{ഗ്രാമം} \ml{ആ} \ml{കവിയെ} \ml{അനുസ്മരിച്ചു}.
\end{quote}

\begin{quote}
\emph{grāmaṁ ā kaviye anusmariccu.}
\end{quote}

\begin{quote}
The village commemorated that poet.
\end{quote}

Past: \ml{അനുസ്മരിച്ചു}, \emph{anusmariccu}.

The person remembered takes -\ml{എ}.

\begin{cousinweb}
\textbf{\ml{അനുസ്മരിക്കുക}} (\emph{anusmarikkuka}) is built on \ml{അനുസ്മരണം} (\emph{anusmaraṇaṁ}), from Sanskrit \emph{anusmaraṇa}, \textquotedblleft{}remembering\textquotedblright{} (Monier-Williams), with the Malayalam verb ending -\ml{ിക്കുക}.
\end{cousinweb}

\begin{grammarlens}[title={What you've built}]
One more word for telling what happened.
\end{grammarlens}

\subsection*{Your turn}
\begin{itemize}
\raggedright
  \item \emph{Say it:} \emph{anusmarikkuka}
  \item \emph{Say it:} \emph{anusmarikkuka}, once more
  \item \emph{Say it:} \emph{grāmaṁ ā kaviye anusmariccu.} — The village commemorated that poet.
  \item \emph{Recall:} say the Malayalam for a statue, then the Malayalam for a stone inscription, then say \emph{anusmarikkuka} again
\end{itemize}

\begin{tcolorbox}[breakable,colback=teal!4,colframe=teal!35!black,title={Before you move on}]
What does \ml{അനുസ്മരിക്കുക} mean? (\textquotedblleft{}To commemorate\textquotedblright{}.) Say it once more.
\end{tcolorbox}
